% Generated by make_paper_figures.py; do not edit values by hand.
\begin{table}[htbp]
\centering
\footnotesize
\setlength{\tabcolsep}{3.5pt}
\begin{tabular*}{\linewidth}{@{\extracolsep{\fill}}lrccc@{}}
\toprule
Input & $n$ & $T_M/T_E$ [interval] & $T_M/T_F$ [interval] & Supported claim \\
\midrule
Conway & 11 & 0.991 [0.448, 3.221] & 1.717 [0.605, 3.925] & None \\
Conway \#2 & 22 & 1.175 [0.582, 4.046] & 1.159 [0.286, 4.324] & None \\
Conway \#3 & 33 & 1.202 [0.292, 2.032] & 0.873 [0.426, 2.014] & None \\
Conway \#8 & 88 & 1.645 [0.474, 4.439] & 0.906 [0.346, 1.539] & None \\
Kinoshita--Terasaka & 11 & 1.374 [0.435, 3.422] & 1.973 [0.305, 3.946] & None \\
\bottomrule
\end{tabular*}
\par\smallskip
\begin{minipage}{\linewidth}\scriptsize
$M$: the normal pipeline with its native generic matching evaluation;
$E,F$: the normal pipeline with exact or factored exact $q=6$.
Nominal pointwise intervals; see the experimental-design assumptions in the text.
All five inputs reached Jones evaluation; the other twelve inputs bypassed it.
No comparison in this seven-round pipeline run met the timing-claim rule.
\end{minipage}
\caption{The five normal-pipeline cases that exercised the selected Jones backend. Timings include input parsing, checked diagram conversion, recognition, and JSON output encoding in a hot process.}
\label{tab:pipeline-selected}
\end{table}
